We now describe the implementation of the protocol given in \cref{fig:protocol_report}.
We instantiate the linearly homomorphic encryption scheme with dense ciphertext space required by $\pdope$ with the Benaloh encryption scheme described in \cref{sec:benaloh_scheme}.
We implement all algorithms in CUDA C$++$ using the NTL \cite{ntl} library.
We additionally make a number of optimizations to the $\chmdss_\np.\reconstruct$ algorithm, which we describe below.
All code can be seen at \repourl{}.

\subsection{Optimizing Aggregation}
\label{sec:star-optimizations}

While the MDSS implementation of $\chmdss$ from \cite{\citeairtag} aggregates efficiently for up to $3150$ shares, it does not scale well for the parameters required by this setting, in which it may be run on sets of $100{,}000$ or more.
Indeed, using the implementation from \cite{\citeairtag}, aggregating from $50{,}000$ shares takes over two hours \emph{per recovered secret}, resulting in a $17$ hour runtime.
Here we first describe the inner-workings of the $\chmdss$ algorithm, and then the optimizations we make to improve the aggregation runtime.

\subsubsection{Algorithm Overview}
\label{sec:algorithm-overview}
Recall from \cref{sec:chmdss} that $\chmdss.\mdssReconstructAlg$ takes a set of $\mdssWindowSize$ shares as input and outputs the secrets corresponding to  values with over $\threshold$ shares.
In our MDSS construction each share is a tuple of field elements $(x, y_1, y_2, \dots, y_\np)$.

The $\chmdss.\mdssReconstructAlg$ algorithm first uses the set of shares to construct a $(c + 1) \times (c + 1)$ matrix, where each entry is a polynomial of max degree $\mdssWindowSize$.
It then puts the matrix into weak Popov form \cite{popov1972invariant}, using the algorithm described in \cite{mulders2003lattice}.
The \cite{mulders2003lattice} algorithm proceeds by iteratively performing elementary row operations (EROs) on the matrix, i.e.\ multiplying a row by a monomial $\polycoeff x^d$ and adding the result to another row.

Once the matrix is in Popov form, $\chmdss.\mdssReconstructAlg$ extracts the shortest vector from the matrix and performs a few simple checks on it.
If the checks pass then the algorithm uses the vector to compute the dealer polynomials $(p_1,\dots,p_{\np})$, from which a secret can be recovered, and if the checks fail the algorithm terminates.
If a secret is recovered $\chmdss.\mdssReconstructAlg$ removes all shares that agree with the recovered polynomials $(p_1,\dots,p_{\np})$ and then repeats, attempting to reconstruct a secret from the remaining $\mdssWindowSize' < \mdssWindowSize$ points.
We refer to a single iteration: building the matrix, computing the Popov form, and attempting to reconstruct the secret, as a \emph{step} of the algorithm.
Each step recovers a single heavy-hitter.
A full description of a step is shown in \cref{fig:chmdss-step}.

\newcommand{\fex}{x}
\newcommand{\fey}{y}
\newcommand{\pvar}{z}
\newcommand{\popov}{\mathsf{popov}}
\newcommand{\mlen}{\Lambda}

\begin{figure}
	\begin{pcvstack}[boxed, center, space=1em]
		\procedure[linenumbering]{$\step\left(\mdssRecThr, \mdssPrivThr, \mdssWindowSize, \set{(\fex_i, \fey_{i,1},\dots,\fey_{i,\np})}_{i \in [\mdssWindowSize]}\right)$}{
			\mlen := \mdssPrivThr + \mdssWindowSize - \mdssRecThr \\
			\pcfor j \in [\np] \pcdo  \\
			\t f_j(\pvar) := \mathsf{LagrInterpol}\left(\set{(\fex_i, \fey_{i, j})}_{i \in [\mdssWindowSize]}\right) \\
			P(\pvar) := \prod_{i = 1}^{\mdssWindowSize}(\pvar - \fex_i) \\
			\text{construct the matrix } M \in \FF[\pvar]^{(\np + 1) \times (\np + 1)} \pcskipln \\
			M := {\begin{pcmbox}\begin{bmatrix}
							z^{d} & f_1(z) & f_2(z) & \ldots & f_n(z) \\
							      & P(z)   &        &        &        \\
							      &        & P(z)   &        &        \\
							      &        &        & \ddots &        \\
							      &        &        &        & P(z)
						\end{bmatrix}\end{pcmbox}}\\
			M_{\popov} := \mathsf{WeakPopovForm}(M)\\
			\pcif M_{\popov} \text{ has exactly one shortest vector } \vec{v} = (v_0, \dots, v_\np) \text{ s.t.\ } |\vec{v}| \leq \mlen \textbf{ and}\pcskipln \\
			v_0 \text{ divides each of } \set{v_1 \cdot \pvar^{\mdssPrivThr}, \dots, v_{\np} \cdot \pvar^{\mdssPrivThr}} \pcthen \\
			\t \pcreturn \left(v_1 \cdot \pvar^{\mdssPrivThr} / v_0, \dots, v_{\np} \cdot \pvar^{\mdssPrivThr}/v_0\right) \\
			\pcelse \pcreturn \bot
		}
	\end{pcvstack}

	\caption{The $\chmdss.\step$ algorithm, adapted from \cite{\citeairtag}}
	\label{fig:chmdss-step}
\end{figure}

